\chapter{Twierdzenie Hahna--Mazurkiewicza i konstrukcja krzywych Peana}
\label{ch:topog26-hahn-mazurkiewicz}

\noindent\textbf{Cel rozdziału.}
W rozdziale~\ref{ch:topog25-continua} wprowadziliśmy
continua Peana i udowodniliśmy istnienie ciągłych
surjekcji odcinka na skończone spójne grafy
oraz na sześciany \([0,1]^n\).
Pozostała ogólna charakteryzacja przestrzeni,
które są ciągłymi obrazami odcinka.
Udowodnimy ją w obu kierunkach,
ze szczególnym naciskiem na konstrukcję
krzywej surjektywnej i uzasadnienie jej ciągłości.
Nie powtarzamy definicji zwartości ani
podstawowych twierdzeń o spójności
z rozdziałów~07--08.

\section{Małe spójne otoczenia w przestrzeni zwartej}
\label{sec:topog26-uniform-connected}

\begin{lemat}[Jednostajna lokalna spójność]
\label{lem:topog26-uniform-local-connected}
Niech \((X,d)\) będzie niepustą zwartą
przestrzenią metryczną, która jest lokalnie spójna.
Dla każdego \(\varepsilon>0\) istnieje
\(\delta>0\) takie, że dowolne
\(a,b\in X\) spełniające \(d(a,b)<\delta\)
należą do wspólnego spójnego
zbioru otwartego \(V\subseteq X\)
o średnicy mniejszej niż \(\varepsilon\).
\end{lemat}
\begin{proof}
Dla każdego \(x\in X\) z lokalnej spójności
wybieramy otwarty spójny zbiór
\(V_x\) taki, że
\(x\in V_x\subseteq B_d(x,\varepsilon/3)\).
Wówczas \(\operatorname{diam}V_x<\varepsilon\).
Z otwartości \(V_x\) wybieramy \(r_x>0\)
takie, że \(B_d(x,2r_x)\subseteq V_x\).
Zwartość \(X\) daje skończone
podpokrycie
\[
 X=\bigcup_{i=1}^{m}B_d(x_i,r_{x_i}).
\]
Niech \(\delta=\min_{1\le i\le m}r_{x_i}>0\).
Jeżeli \(d(a,b)<\delta\), wybieramy
\(i\) z \(a\in B_d(x_i,r_{x_i})\).
Wtedy
\[
 d(b,x_i)\le d(b,a)+d(a,x_i)
       <\delta+r_{x_i}\le2r_{x_i}.
\]
Oba punkty należą zatem do \(V_{x_i}\),
który jest otwarty, spójny
i ma średnicę mniejszą niż \(\varepsilon\).
\end{proof}

\begin{lemat}[Łańcuchy w zbiorach spójnych]
\label{lem:topog26-chain}
Niech \((Y,d)\) będzie spójną przestrzenią metryczną
i \(\delta>0\). Dla dowolnych \(a,b\in Y\)
istnieją \(m\ge1\) i punkty
\(a=y_0,y_1,\ldots,y_m=b\)
takie, że \(d(y_{j-1},y_j)<\delta\)
dla wszystkich \(j=1,\ldots,m\).
\end{lemat}
\begin{proof}
Ustalmy \(a\in Y\) i oznaczmy przez \(A\)
zbiór wszystkich punktów, które można
połączyć z \(a\) skończonym łańcuchem
o odległościach kolejnych punktów
mniejszych od \(\delta\).
Zbiór \(A\) jest niepusty, gdyż zawiera \(a\).
Jeżeli \(y\in A\), to \(B_d(y,\delta)\subseteq A\):
wystarczy dołączyć kolejny punkt do łańcucha.
Zatem \(A\) jest otwarty.
Jeżeli \(z\in Y\setminus A\), to kula
\(B_d(z,\delta)\) nie spotyka \(A\),
bo w przeciwnym razie \(z\in A\).
Dopełnienie \(A\) jest więc także otwarte.
Ze spójności \(Y\) otrzymujemy \(A=Y\).
\end{proof}

\section{Budowa drogi z coraz drobniejszych łańcuchów}
\label{sec:topog26-paths}

\begin{twierdzenie}[Małe drogi w continuach Peana]
\label{thm:topog26-small-paths}
Niech \((X,d)\) będzie niepustym
continuum metrycznym, lokalnie spójnym.
Dla każdego \(\eta>0\) istnieje
\(\delta>0\) takie, że dla dowolnych
\(a,b\in X\) spełniających \(d(a,b)<\delta\)
istnieje droga \(\gamma\colon[0,1]\to X\)
łącząca \(a\) z \(b\), której obraz
ma średnicę mniejszą niż \(\eta\).
W szczególności \(X\) jest
lokalnie spójna drogowo i spójna drogowo.
\end{twierdzenie}
\begin{proof}
Ustalmy \(\eta>0\).
Wybierzmy dodatni ciąg \((\varepsilon_n)_{n\ge0}\)
z
\(\sum_{n=0}^{\infty}\varepsilon_n<\eta/4\),
na przykład
\(\varepsilon_n=\eta\,2^{-n-4}\).
Z Lematu~\ref{lem:topog26-uniform-local-connected}
dla każdego \(n\) wybieramy
\(\delta_n>0\) takie, że
każde dwa punkty odległe o mniej
niż \(\delta_n\) należą do
wspólnego spójnego zbioru
o średnicy mniejszej niż \(\varepsilon_n\).
Możemy dodatkowo żądać
\(\delta_{n+1}<\delta_n/2\).
Niech \(\delta=\delta_0\).

\emph{Krok 1: konstruowanie wartości na gęstym zbiorze.}
Dla \(a,b\in X\) z \(d(a,b)<\delta_0\)
przypiszmy wartości \(a,b\)
końcom odcinka \([0,1]\).
Załóżmy indukcyjnie, że skonstruowaliśmy
skończony podział
\[
 0=t_{n,0}<t_{n,1}<\cdots<t_{n,N_n}=1
\]
oraz wartości \(x_{n,j}\in X\)
w jego wierzchołkach, zachowujące
\(x_{n,0}=a\), \(x_{n,N_n}=b\),
takie że
\[
 d(x_{n,j-1},x_{n,j})<\delta_n.
 \tag{1}
\]
Dla \(n=0\) przyjmujemy
podział o jednym odcinku i
wartości \(a,b\).

Dla każdego odcinka
\([t_{n,j-1},t_{n,j}]\)
Lemat~\ref{lem:topog26-uniform-local-connected}
dostarcza spójnego zbioru
\(C_{n,j}\subseteq X\),
zawierającego oba końcowe punkty
i mającego średnicę mniejszą
niż \(\varepsilon_n\).
Lemat~\ref{lem:topog26-chain}
pozwala połączyć te punkty
skończonym \(\delta_{n+1}\)-łańcuchem
złożonym z punktów \(C_{n,j}\).
Dzielimy rozpatrywany odcinek
czasu na tyle równych części,
ile łańcuch ma odcinków,
i przypisujemy nowym wierzchołkom
kolejne punkty łańcucha.
Jeżeli wybrany łańcuch ma
tylko jeden odcinek, dopisujemy
na jego początku jeszcze raz
pierwszy punkt. Dzięki temu
każdy odcinek czasu zostaje
podzielony na co najmniej dwie części.
Na wspólnych końcach starych
odcinków zachowujemy poprzednio
przypisane wartości.

W ten sposób otrzymujemy
podział poziomu \(n+1\),
zgodny z poprzednim,
a na każdym jego odcinku
kolejne wartości spełniają (1)
z indeksem \(n+1\).
Długość najdłuższego odcinka
podziału poziomu \(n\)
jest nie większa niż \(2^{-n}\).

Niech \(D\subseteq[0,1]\) będzie
sumą zbiorów wierzchołków
wszystkich podziałów.
Jest to zbiór gęsty:
średnice odcinków podziału dążą do zera.
Ponieważ wartości na starych wierzchołkach
nie są zmieniane, konstrukcja
wyznacza jednoznaczne odwzorowanie
\(\gamma_0\colon D\to X\).

\emph{Krok 2: oszacowanie oscylacji.}
Niech
\(S_n=\sum_{k=n}^{\infty}\varepsilon_k\).
Ustalmy dowolny odcinek \(J\)
podziału na poziomie \(n\).
Wszystkie wartości przypisane
jego wierzchołkom przy kolejnych
podziałach są odległe o mniej
niż \(S_n\) od wartości
na lewym końcu \(J\):
na pierwszym etapie nowy punkt
należy do \(C_{n,j}\),
więc odległość jest mniejsza
niż \(\varepsilon_n\).
W kolejnych etapach
nowo dodany punkt leży w
spójnym zbiorze o średnicy
mniejszej niż odpowiednie
\(\varepsilon_k\), zawierającym
końce jego odcinka rodzicielskiego.
Indukcja daje sumę oszacowań
\(\varepsilon_n+\cdots+\varepsilon_{N-1}\)
na dowolnym poziomie \(N>n\).

W konsekwencji oscylacja \(\gamma_0\)
na \(D\cap J\) jest co najwyżej \(2S_n\).
Ponieważ wartości przypisane
prawemu końcowi \(J\)
też należą do tego zbioru,
każdy punkt \(D\cap J\)
jest od niego odległy
o co najwyżej \(2S_n\).
Dla dwóch sąsiednich odcinków
poziomu \(n\) dostajemy więc
oszacowanie \(4S_n\),
wykorzystując ich wspólny wierzchołek.

Teraz dla dowolnego \(\rho>0\)
wybieramy \(n\) tak, aby \(4S_n<\rho\).
Ponieważ podział poziomu \(n\)
jest skończony, ma dodatnią
najmniejszą długość odcinka \(h_n\).
Jeżeli \(u,v\in D\) i
\(|u-v|<h_n\), to \(u,v\)
należą do tego samego lub
do dwóch sąsiednich odcinków
podziału poziomu \(n\).
Zatem \(d(\gamma_0(u),\gamma_0(v))<\rho\).
Odwzorowanie \(\gamma_0\)
jest jednostajnie ciągłe na \(D\).

\emph{Krok 3: przedłużenie do drogi.}
Zwartość metrycznego \(X\)
implikuje jego zupełność.
Dla każdego \(t\in[0,1]\)
wybieramy ciąg \(t_k\in D\)
z \(t_k\to t\).
Jednostajna ciągłość \(\gamma_0\)
sprawia, że \((\gamma_0(t_k))\)
jest ciągiem Cauchy'ego,
więc ma granicę w \(X\).
Granica nie zależy od wyboru
ciągu \(t_k\), ponieważ dla
dwóch takich ciągów
ich odległość w parametrze
dąży do zera.
Definiujemy \(\gamma(t)\)
jako tę granicę.
Otrzymujemy ciągłe przedłużenie
\(\gamma\colon[0,1]\to X\)
z \(\gamma(0)=a,\gamma(1)=b\).

Stosując oszacowanie z kroku 2
do całego początkowego odcinka
i przechodząc do granicy,
otrzymujemy
\(d(\gamma(t),a)\le S_0\)
dla wszystkich \(t\).
Wobec tego
\(\operatorname{diam}\gamma([0,1])
\le2S_0<\eta\).
Udowodniliśmy główną tezę.

\emph{Krok 4: spójność drogowa lokalna i globalna.}
Niech \(U\subseteq X\) będzie otwarty,
a \(x\in U\).
Wybierzmy \(r>0\)
z \(B_d(x,r)\subseteq U\).
Konstrukcja powyżej z \(\eta<r\)
daje \(\delta_x>0\) takie, że
każdy \(y\in B_d(x,\delta_x)\)
można połączyć z \(x\)
drogą o obrazie zawartym
w \(B_d(x,r)\subseteq U\).
Dlatego składowa drogowa
\(U\) zawierająca \(x\)
zawiera otwarte otoczenie \(x\).
Stosując ten argument do
każdego punktu tej składowej,
widzimy, że jest ona otwarta.
Wszystkie składowe drogowe
wszystkich otwartych podzbiorów
są więc otwarte.
Z charakterystyki
\ref{thm:topog-spojnosc-lokalna-komponenty}
otrzymujemy lokalną spójność drogową.
Ponieważ \(X\) jest spójna,
wniosek~\ref{cor:topog-spojnosc-lokalna-drogowa}
daje spójność drogową \(X\).
\end{proof}

\begin{uwaga}[Dlaczego nie wystarczy spójność lokalna bez zwartości?]
\label{rem:topog26-compactness}
Zwartość była niezbędna do uzyskania
jednego \(\delta\) działającego
dla wszystkich punktów \(X\).
Ponadto jej konsekwencja, zupełność,
pozwoliła przedłużyć funkcję
z gęstego zbioru parametrów.
Sama spójność lokalna w dowolnej
przestrzeni metrycznej nie uzasadnia
powyższej konstrukcji.
\end{uwaga}

\section{Dołączanie krótkich wycieczek do drogi}
\label{sec:topog26-excursions}

\begin{lemat}[Uzupełnienie obrazu drogi o skończony zbiór]
\label{lem:topog26-excursion}
Niech \((X,d)\) będzie continuum Peana,
\(h\colon[0,1]\to X\) drogą,
\(\varepsilon>0\) i \(\delta>0\)
takimi, że każdą parę punktów
odległych o mniej niż \(\delta\)
można połączyć drogą o obrazie
średnicy mniejszej niż \(\varepsilon\).
Niech \(E\subseteq X\)
będzie niepustym skończonym zbiorem,
spełniającym
\[
 \forall y\in E\quad
 \exists t_y\in[0,1]\quad
 d(y,h(t_y))<\delta.
\]
Wówczas istnieje droga
\(g\colon[0,1]\to X\)
taka, że
\[
 g(0)=h(0),\qquad g(1)=h(1),\qquad
 E\subseteq g([0,1]),\qquad
 \sup_{t\in[0,1]}d(g(t),h(t))<2\varepsilon.
\]
\end{lemat}
\begin{proof}
Zbierzmy wszystkie różne wartości
\(t_y\) w skończony zbiór
\(T=\{s_1,\ldots,s_m\}\subseteq[0,1]\).
Dla każdego \(s_i\) wybieramy
domknięty przedział
\(J_i=[a_i,b_i]\subseteq[0,1]\)
z \(a_i<b_i\) i \(s_i\in J_i\),
tak aby przedziały \(J_i\)
były parami rozłączne oraz
\[
 \operatorname{diam}h(J_i)<\varepsilon.
 \tag{2}
\]
Można to zrobić dzięki jednostajnej
ciągłości \(h\) oraz skończoności \(T\).
Gdy \(s_i=0\) lub \(1\),
wybieramy przedział jednostronny.

Dla każdego \(y\in E\) z \(t_y=s_i\)
Twierdzenie~\ref{thm:topog26-small-paths}
daje drogę
\(\alpha_y\colon[0,1]\to X\)
od \(h(s_i)\) do \(y\),
o obrazie średnicy mniejszej
niż \(\varepsilon\).
Połączenie \(\alpha_y\)
z drogą przebieganą w przeciwnym kierunku
jest pętlą zaczynającą się
i kończącą w \(h(s_i)\),
której obraz zawiera \(y\).

Na każdym \(J_i\) budujemy
odcinkami ciągłą drogę
przebiegającą kolejno:
odcinek starej drogi \(h\)
od \(a_i\) do \(s_i\),
wszystkie pętle do punktów
\(y\) z \(t_y=s_i\),
a następnie odcinek \(h\)
od \(s_i\) do \(b_i\).
Parametryzujemy te finitnie
liczne fragmenty liniowo
na kolejnych częściach \(J_i\).
Jeśli \(s_i=a_i\) lub \(s_i=b_i\),
odpowiedni fragment stały można pominąć.
Na dopełnieniu sumy przedziałów \(J_i\)
kładziemy \(g(t)=h(t)\).
Nowe i stare fragmenty zgadzają się
w punktach \(a_i,b_i\),
zatem \(g\) jest ciągła.
Ma wymagane końce i zawiera każdy punkt \(E\).

Dla \(t\in J_i\) punkt \(g(t)\)
jest albo punktem z \(h(J_i)\),
albo punktem pętli do któregoś \(y\).
W obu przypadkach
\(d(g(t),h(s_i))<\varepsilon\),
na mocy (2) i średnicy obrazu
każdej pętli.
Ponieważ również
\(d(h(t),h(s_i))<\varepsilon\),
mamy \(d(g(t),h(t))<2\varepsilon\).
Poza przedziałami \(J_i\)
funkcje są identyczne.
To kończy dowód.
\end{proof}

\begin{figure}[htbp]
\centering
\begin{tikzpicture}[x=1.1cm,y=1.05cm,>=stealth]
  \draw[gray!60,thick] (0,0) -- (1,0.35) -- (2,0.15)
    -- (3,0.5) -- (4,0.2) -- (5,0.7) -- (6,0.45);
  \draw[very thick,blue!70!black]
    (0,0) -- (1,0.35) -- (2,0.15) -- (2.6,0.42)
    .. controls (2.75,1.55) and (3.4,1.35) .. (3.0,0.5)
    .. controls (3.3,1.45) and (3.75,1.05) .. (3.0,0.5)
    -- (4,0.2) -- (5,0.7) -- (6,0.45);
  \fill[blue!70!black] (3,0.5) circle (2pt);
  \fill[blue!70!black] (3.2,1.13) circle (2pt);
  \node[font=\small,anchor=west] at (3.37,1.40)
     {punkt dołączony krótką pętlą};
  \node[font=\small,anchor=west] at (.1,-.40)
     {szara: droga wyjściowa};
  \node[font=\small,blue!70!black,anchor=west] at (3.6,-.40)
     {niebieska: droga po poprawce};
\end{tikzpicture}
\caption{Schemat lokalnej poprawki drogi:
na krótkim przedziale parametru
dołączamy pętlę przechodzącą przez
nowy punkt, zachowując końce całej drogi.
Rysunek ma charakter ilustracyjny.}
\label{fig:topog26-excursion}
\end{figure}

\section{Twierdzenie Hahna--Mazurkiewicza: wystarczalność}
\label{sec:topog26-sufficiency}

\begin{twierdzenie}[Twierdzenie Hahna--Mazurkiewicza, kierunek konstrukcyjny]
\label{thm:topog26-hm-sufficient}
Niech \(X\) będzie niepustym
continuum metrycznym, lokalnie spójnym.
Istnieje ciągła surjekcja
\(f\colon[0,1]\to X\).
\end{twierdzenie}
\begin{proof}
Ustalmy zgodną metrykę \(d\) na \(X\).
Wybierzmy ciąg dodatni
\(\varepsilon_n=2^{-n-3}\)
oraz, na mocy
Twierdzenia~\ref{thm:topog26-small-paths},
ciąg \((\delta_n)_{n\ge0}\) liczb dodatnich
taki, że
\(\delta_{n+1}<\delta_n/2\)
i każdą parę punktów
odległych o mniej niż \(\delta_n\)
można połączyć drogą
o obrazie średnicy mniejszej
niż \(\varepsilon_n\).

\emph{Krok 1: pierwsza droga o obrazie gęstym
w skali \(\delta_0\).}
Zwartość \(X\) daje
skończony zbiór
\(E_0=\{y_1,\ldots,y_m\}\subseteq X\)
taki, że dla każdego \(x\in X\)
istnieje \(y_j\in E_0\)
z \(d(x,y_j)<\delta_0/2\).
Na mocy
Twierdzenia~\ref{thm:topog26-small-paths}
przestrzeń \(X\) jest spójna drogowo.
Możemy więc połączyć kolejno
punkty \(y_1,\ldots,y_m\)
drogami i złożyć je w jedną
ciągłą drogę
\(f_0\colon[0,1]\to X\),
której obraz zawiera \(E_0\).
Gdy \(m=1\), bierzemy
drogę stałą.
Obraz \(f_0([0,1])\)
jest \(\delta_0/2\)-gęsty w \(X\).

\emph{Krok 2: poprawianie drogi.}
Załóżmy, że skonstruowaliśmy
drogę \(f_n\colon[0,1]\to X\),
której obraz jest
\(\delta_n/2\)-gęsty.
Zwartość \(X\) pozwala
wybrać skończony zbiór
\(E_{n+1}\subseteq X\)
będący \(\delta_{n+1}/2\)-siecią,
tzn.
\[
 \forall x\in X\quad
 \exists y\in E_{n+1}\quad
 d(x,y)<\delta_{n+1}/2.
\]
Z gęstości obrazu \(f_n\)
wynika, że dla każdego
\(y\in E_{n+1}\) istnieje
\(t_y\in[0,1]\)
z \(d(y,f_n(t_y))<\delta_n/2<\delta_n\).
Lemat~\ref{lem:topog26-excursion},
z parametrami \(\delta_n\)
i \(\varepsilon_n\),
daje nową drogę
\(f_{n+1}\colon[0,1]\to X\),
której obraz zawiera
\(E_{n+1}\) oraz
\[
 \sup_{t\in[0,1]}
 d(f_{n+1}(t),f_n(t))
 \le2\varepsilon_n.
 \tag{3}
\]
Wobec tego obraz \(f_{n+1}\)
jest \(\delta_{n+1}/2\)-gęsty.
Rekurencja konstruuje
drogi \(f_n\) dla wszystkich \(n\).

\emph{Krok 3: istnienie granicy jednostajnej.}
Dla \(m>n\), stosując nierówność
trójkąta oraz (3), mamy
\[
 \sup_{t\in[0,1]}d(f_m(t),f_n(t))
 \le \sum_{k=n}^{m-1}2\varepsilon_k.
\]
Szereg \(\sum_{k\ge0}2\varepsilon_k\)
jest zbieżny, więc \((f_n)\)
jest ciągiem Cauchy'ego
w metryce zbieżności jednostajnej.
Ponieważ \(X\) jest zupełna,
dla każdego \(t\) istnieje
granica \(f(t)=\lim_n f_n(t)\),
a powyższe oszacowanie
daje zbieżność jednostajną.
Funkcja \(f\colon[0,1]\to X\)
jest ciągła jako granica
jednostajna funkcji ciągłych.

\emph{Krok 4: surjektywność granicy.}
Niech \(x\in X\).
Dla każdego \(n\) istnieje
\(t_n\in[0,1]\) z
\(d(x,f_n(t_n))<\delta_n/2\).
Zatem
\[
 d(x,f(t_n))
 \le \frac{\delta_n}{2}
  +\sup_{t\in[0,1]}d(f_n(t),f(t))
 \le \frac{\delta_n}{2}
  +\sum_{k=n}^{\infty}2\varepsilon_k
 \longrightarrow0.
\]
Ponieważ \(f([0,1])\)
jest zwartym, a więc domkniętym
podzbiorem Hausdorffowego \(X\),
istnienie punktów \(f(t_n)\)
zbiegających do \(x\)
implikuje \(x\in f([0,1])\).
Punkt \(x\) był dowolny,
więc \(f([0,1])=X\).
\end{proof}

\begin{uwaga}[Charakter konstrukcyjny]
\label{rem:topog26-constructive}
Powyższy dowód nie zakłada
istnienia krzywej wypełniającej \(X\)
na żadnym wcześniejszym etapie.
Poziom \(n\) dostarcza drogi
o obrazie gęstym w skali
\(\delta_n/2\), a koszt poprawki
jest kontrolowany przez
\(2\varepsilon_n\).
Suma kosztów jest skończona,
dlatego granica jest ciągła;
gęstość w coraz drobniejszej
skali i domkniętość obrazu
zapewniają surjektywność.
\end{uwaga}

\section{Własności konieczne obrazów odcinka}
\label{sec:topog26-necessity}

\begin{lemat}[Obraz odcinka spełnia drugi aksjomat przeliczalności]
\label{lem:topog26-image-second-countable}
Niech \(X\) będzie przestrzenią
Hausdorffa, a
\(f\colon[0,1]\to X\)
ciągłą surjekcją.
Wówczas \(X\) ma
przeliczalną bazę topologii.
\end{lemat}
\begin{proof}
Niech \(\mathcal B\)
będzie przeliczalną bazą
topologii odcinka \([0,1]\),
na przykład rodziną
przecięć przedziałów
o wymiernych końcach z \([0,1]\).
Rodzina \(\mathcal V\)
wszystkich skończonych
sum zbiorów z \(\mathcal B\)
jest przeliczalna.
Dla \(V\in\mathcal V\)
definiujemy
\[
 W_V=X\setminus f([0,1]\setminus V).
\]
Zbiór \([0,1]\setminus V\)
jest domknięty w zwartym odcinku,
a więc zwarty.
Jego obraz przez \(f\)
jest zwarty, czyli domknięty
w Hausdorffowym \(X\).
Zatem \(W_V\) jest otwarty.

Pokażemy, że
\(\{W_V:V\in\mathcal V\}\)
jest bazą \(X\).
Niech \(x\in O\), gdzie
\(O\subseteq X\) jest otwarty.
Zbiór \(f^{-1}(\{x\})\)
jest niepusty i zwarty:
jest domknięty w \([0,1]\),
ponieważ \(X\) jest Hausdorffa.
Jest zawarty w otwartym zbiorze
\(f^{-1}(O)\).
Dla każdego jego punktu
dobieramy \(B\in\mathcal B\)
zawierający ten punkt
i spełniający \(B\subseteq f^{-1}(O)\).
Ze zwartości włókna
wybieramy skończenie wiele
takich elementów bazy,
których suma \(V\in\mathcal V\)
spełnia
\[
 f^{-1}(\{x\})\subseteq V
 \subseteq f^{-1}(O).
\]
Pierwsza inkluzja daje
\(x\notin f([0,1]\setminus V)\),
a zatem \(x\in W_V\).
Druga daje \(W_V\subseteq O\):
jeśli \(y\notin O\),
każdy jego przeciwobraz
leży poza \(V\), więc
\(y\in f([0,1]\setminus V)\).
Otrzymaliśmy \(x\in W_V\subseteq O\).
Wszystkie zbiory bazy
tworzą rodzinę przeliczalną.
\end{proof}

\begin{lemat}[Ciągły obraz odcinka jest lokalnie spójny]
\label{lem:topog26-image-locally-connected}
Niech \(X\) będzie przestrzenią
Hausdorffa, a
\(f\colon[0,1]\to X\)
ciągłą surjekcją.
Wówczas \(X\) jest lokalnie spójna.
\end{lemat}
\begin{proof}
Ustalmy otwarty \(O\subseteq X\)
i punkt \(x\in O\).
Włókno
\(F=f^{-1}(\{x\})\)
jest niepuste i zwarte
oraz leży w \(f^{-1}(O)\).
Dla każdego \(t\in F\)
istnieje domknięty przedział
\(I_t\subseteq[0,1]\)
taki, że \(t\) leży w jego
wnętrzu względnym,
\(I_t\subseteq f^{-1}(O)\)
oraz \(I_t\cap F\neq\varnothing\).
Dla \(t=0\) lub \(1\)
używamy przedziałów jednostronnych.
Ze zwartości \(F\)
wybieramy skończoną
rodzinę \(I_1,\ldots,I_m\),
których wnętrza względne
pokrywają \(F\).
Niech
\[
 V=\bigcup_{j=1}^{m}\operatorname{int}_{[0,1]}I_j,
 \qquad K=\bigcup_{j=1}^{m}f(I_j).
\]
Każdy \(f(I_j)\) jest
spójny jako ciągły obraz
przedziału, zawiera \(x\),
gdyż \(I_j\) zawiera
odpowiedni punkt włókna \(F\),
i jest zawarty w \(O\).
Suma skończenie wielu
spójnych zbiorów mających
wspólny punkt \(x\)
jest spójna. Zatem
\(K\subseteq O\)
jest spójny i zawiera \(x\).

Zdefiniujmy
\[
 W=X\setminus f([0,1]\setminus V).
\]
Zbiór \(W\) jest otwarty,
ponieważ \([0,1]\setminus V\)
jest zwarty, a \(X\)
jest Hausdorffa.
Ponieważ \(F\subseteq V\),
mamy \(x\in W\).
Z kolei \(W\subseteq K\):
dla \(y\in W\) surjektywność
\(f\) daje \(t\) z \(f(t)=y\);
z definicji \(W\) wszystkie
takie punkty leżą w \(V\),
a więc należą do pewnego \(I_j\).
Otrzymaliśmy otwarte
otoczenie \(W\ni x\)
z \(W\subseteq K\subseteq O\),
gdzie \(K\) jest spójny.

Niech \(C\) będzie składową
spójną \(O\) zawierającą \(x\).
Ponieważ \(K\) jest spójny
i \(x\in K\cap C\),
mamy \(K\subseteq C\)
na mocy maksymalności składowej.
Zatem \(W\subseteq C\).
Powtarzając to rozumowanie
dla każdego punktu \(C\)
i tego samego \(O\),
widzimy, że \(C\)
jest otwarta w \(X\).
Każda składowa spójna
dowolnego otwartego \(O\)
jest więc otwarta.
Twierdzenie~\ref{thm:topog-spojnosc-lokalna-komponenty}
daje lokalną spójność \(X\).
\end{proof}

\begin{twierdzenie}[Twierdzenie Hahna--Mazurkiewicza]
\label{thm:topog26-hahn-mazurkiewicz}
Niech \(X\) będzie niepustą
przestrzenią topologiczną Hausdorffa.
Następujące warunki są równoważne:
\begin{enumerate}[label=\textup{(\roman*)}]
\item istnieje ciągła surjekcja
\(f\colon[0,1]\to X\);
\item \(X\) jest zwarta, spójna,
lokalnie spójna i metryzowalna.
\end{enumerate}
\end{twierdzenie}
\begin{proof}
Jeżeli zachodzi \textup{(ii)},
to \(X\) jest continuum Peana,
a implikację \textup{(ii)}\(\Rightarrow\)\textup{(i)}
udowodniliśmy konstrukcyjnie
w Twierdzeniu~\ref{thm:topog26-hm-sufficient}.

Załóżmy \textup{(i)}.
Jako ciągły obraz
niepustego zwartego
i spójnego odcinka
\(X\) jest niepusta,
zwarta i spójna.
Ponadto
Lemat~\ref{lem:topog26-image-locally-connected}
daje lokalną spójność \(X\).
Na mocy
Lematu~\ref{lem:topog26-image-second-countable}
przestrzeń \(X\)
spełnia drugi aksjomat
przeliczalności.
Każda zwarta przestrzeń
Hausdorffa jest regularna
i spełnia \(T_1\);
zatem twierdzenie Urysohna
o metryzowalności
\ref{thm:topog-urysohn-metryzowalnosc}
zapewnia metryzowalność \(X\).
To dowodzi \textup{(i)}\(\Rightarrow\)\textup{(ii)}.
\end{proof}

\section{Przykłady, kontrprzykłady i konsekwencje}
\label{sec:topog26-examples}

\begin{wniosek}[Continua Peana są spójne drogowo]
\label{cor:topog26-peano-path-connected}
Niech \(X\) będzie continuum metrycznym
i lokalnie spójnym.
Wówczas \(X\) jest
spójna drogowo oraz lokalnie spójna drogowo.
\end{wniosek}
\begin{proof}
Jest to ostatnia część
Twierdzenia~\ref{thm:topog26-small-paths}.
Można także wyprowadzić
samą spójność drogową
z Twierdzenia~\ref{thm:topog26-hahn-mazurkiewicz},
ponieważ ciągły obraz
drogi jest spójny drogowo,
ale nie dawałoby to
jednostajnego oszacowania
średnicy małych dróg.
\end{proof}

\begin{przyklad}[Zwartość i spójność nie wystarczą]
\label{ex:topog26-solenoid}
Solenoid dyadyczny \(\Sigma_2\),
skonstruowany w
Definicji~\ref{def:topog25-solenoid},
jest zwartym, spójnym,
metryzowalnym continuum.
Nie jest jednak lokalnie spójny
na mocy
Twierdzenia~\ref{thm:topog25-solenoid-not-locally-connected}.
Z Twierdzenia~\ref{thm:topog26-hahn-mazurkiewicz}
wynika, że nie istnieje
ciągła surjekcja
\([0,1]\to\Sigma_2\).
Jest to przykład pokazujący
istotność założenia
lokalnej spójności.
\end{przyklad}

\begin{przyklad}[Niespójność jest osobną przeszkodą]
\label{ex:topog26-cantor}
Dyskretny zbiór \(D=\{0,1\}\)
jest zwarty, Hausdorffa, metryzowalny
i lokalnie spójny: każdy singleton
jest otwartym spójnym otoczeniem
swojego punktu. Nie jest jednak
spójny, więc nie może być ciągłym
obrazem całego odcinka \([0,1]\).

Klasyczny zbiór Cantora jest również
zwarty i metryzowalny, ale nie jest
ani spójny, ani lokalnie spójny:
każdy jego niepusty spójny
podzbiór jest singletonem,
a żaden singleton nie jest otwarty.
Można z niego natomiast konstruować
ciągłe surjekcje na sześciany,
zgodnie z
Twierdzeniem~\ref{thm:topog25-space-filling}.
\end{przyklad}

\begin{przyklad}[Okrąg, grafy i sześciany]
\label{ex:topog26-classical}
Okrąg \(S^1\), każdy skończony
spójny graf topologiczny
oraz każdy sześcian \([0,1]^n\)
są continuami Peana.
Dlatego są ciągłymi
obrazami odcinka.
Dla grafów i sześcianów
konstrukcje surjekcji
zostały już podane w
Twierdzeniach
\ref{thm:topog25-finite-graph-peano}
i~\ref{thm:topog25-space-filling};
obecne twierdzenie
obejmuje je jako przypadki
ogólniejszej charakteryzacji.
\end{przyklad}

\begin{uwaga}[Istotność założenia Hausdorffa]
\label{rem:topog26-metrizability}
Bez aksjomatu Hausdorffa
ciągły obraz odcinka
nie musi być metryzowalny.
Istotnie, weźmy
\(Y=\{0,1\}\) z topologią
\(\{\varnothing,\{1\},Y\}\)
i określmy
\(f\colon[0,1]\to Y\) wzorem
\(f(0)=0\), \(f(t)=1\) dla \(t>0\).
Odwzorowanie jest surjektywne i ciągłe,
gdyż przeciwobrazem jedynego
nietrywialnego zbioru otwartego
jest \((0,1]\), otwarty w \([0,1]\).
Jednak \(Y\) nie jest Hausdorffa,
a zatem nie jest metryzowalna.
To uzasadnia jawne założenie
Hausdorffowości w
Twierdzeniu~\ref{thm:topog26-hahn-mazurkiewicz}.
\end{uwaga}

\par\smallskip\noindent\textbf{Podsumowanie.}
Twierdzenie Hahna--Mazurkiewicza
charakteryzuje ciągłe obrazy
odcinka wśród niepustych
przestrzeni Hausdorffa.
Konstrukcja wystarczalności
ma dwa niezależne etapy:
otrzymanie małych dróg
z lokalnej spójności
i zwartości oraz dołączanie
coraz drobniejszych
wycieczek do skończonych sieci.
Jednostajna zbieżność
poprawek daje krzywą ciągłą,
a coraz większa gęstość
jej przybliżeń
gwarantuje surjektywność.
W kierunku koniecznym
zwartość włókien
ciągłego odwzorowania odcinka
daje lokalną spójność
i drugi aksjomat przeliczalności,
a twierdzenie Urysohna
daje metryzowalność.